% ============================================================
% Amazon-VT CFP 2026-2027 — 1-Page Abstract
% Provenance-Aware Agentic Knowledge Graph Systems
% ============================================================

\documentclass[10pt, letterpaper]{article}

% --- Packages ---
\usepackage[utf8]{inputenc}
\usepackage[T1]{fontenc}
\usepackage{lmodern}
\usepackage[margin=0.85in, top=0.75in, bottom=0.75in]{geometry}
\usepackage{hyperref}
\usepackage{natbib}
\usepackage{microtype}
\usepackage{parskip}
\setlength{\parskip}{0.35em}

% Tighten spacing to fit 1 page
\usepackage{titlesec}
\titlespacing*{\section}{0pt}{0.8em}{0.4em}

% --- Metadata ---
\title{\vspace{-2.5em}Provenance-Aware Agentic Knowledge Graph Systems\\for Reliable Multi-Source Grounded AI\vspace{-0.5em}}
\author{}
\date{}

% ============================================================
\begin{document}

\maketitle
\thispagestyle{empty}

\vspace{-2.5em}

\noindent\textbf{Primary Topic Area:} Knowledge Grounding \hfill
\textbf{Secondary Topic Area:} Agentic AI

\vspace{0.3em}

Large Language Models (LLMs) and generative AI systems increasingly rely on retrieval-based mechanisms to incorporate external knowledge during inference. Most current retrieval-augmented generation (RAG) systems depend primarily on vector similarity search over unstructured document collections. While this approach improves relevance, it provides limited support for structured reasoning, provenance tracking, and verifiable evidence chains. As agentic AI systems become capable of autonomous task execution, planning, and decision-making, the absence of grounded, traceable knowledge sources introduces risks including hallucination, unverifiable outputs, and brittle reasoning across heterogeneous information sources.

This project proposes \textbf{Provenance-Aware Agentic Knowledge Graph Systems (PA-AKG)}, a hybrid architecture that integrates knowledge graphs, multi-source retrieval, and agent orchestration to support reliable grounded reasoning. The key insight of this work is that structured knowledge representations can serve as a reasoning substrate for generative systems, enabling agents to combine symbolic relationships, unstructured text retrieval, and dynamic reasoning workflows while maintaining explicit provenance and evidence tracking. In this framework, agents interact with both vector-based retrieval systems and structured knowledge graphs, dynamically selecting retrieval pathways based on task context and reasoning requirements.

The proposed research will develop a unified architecture for \textbf{multi-source knowledge grounding across heterogeneous data environments}, including knowledge graphs, relational databases, scientific literature, and streaming or domain-specific datasets. Agent workflows will be designed to support reasoning chains that explicitly reference structured entities, relationships, and supporting evidence. A key component of the system is an automated \textbf{provenance generation and validation mechanism}, enabling agents to produce traceable reasoning paths that link generated outputs to supporting knowledge sources. This capability allows both automated and human verification of system outputs, improving reliability in complex reasoning tasks.

The project will investigate several core research questions: (1)~How can heterogeneous knowledge sources---including knowledge graphs, databases, and document corpora---be integrated into a unified grounding framework for generative systems? (2)~How should agentic workflows dynamically select between structured and unstructured retrieval pathways during reasoning tasks? (3)~What mechanisms enable scalable provenance generation and verification for agent-generated outputs? (4)~How does structured knowledge grounding impact reasoning accuracy, robustness, and computational efficiency compared to conventional RAG systems?

To answer these questions, we will develop a \textbf{hybrid KG-RAG architecture}, implement agent orchestration workflows that integrate structured reasoning with generative models, and design evaluation benchmarks that measure grounded reasoning quality, provenance completeness, and system efficiency. Experimental evaluation will compare three system paradigms: traditional vector-based RAG, hybrid structured retrieval systems, and fully knowledge-graph-driven agentic reasoning architectures.

The research will be evaluated across representative application domains that require reliable knowledge grounding and complex reasoning, including scientific research discovery and engineering decision support systems. These domains provide challenging scenarios where generative models must integrate structured constraints, heterogeneous knowledge sources, and multi-step reasoning processes.

Expected outcomes include: a novel architecture for \textbf{provenance-aware knowledge-grounded agentic AI systems}; methods for integrating structured knowledge graphs with generative retrieval pipelines; evaluation benchmarks for grounded reasoning and provenance validation in agentic systems; and an open reference implementation enabling reproducible research and future extensions.

By enabling generative AI systems to reason over structured knowledge while maintaining explicit evidence chains, this work aims to significantly improve the \textbf{reliability, transparency, and robustness of agentic AI systems} operating over complex real-world knowledge environments.

\end{document}
